\begin{titlepage}
    \centering
    \setstretch{1.0} % Title page usually uses single spacing for grouping

    % 1. Institution Info (12pt)
    {\fontsize{12pt}{14pt}\selectfont 
    TARTU ÜLIKOOLI \\
    Sotsiaalteaduste Valdkond \\
    Narva Kolledž \par}

    \vspace{2cm}

    % 2. Curriculum (12pt)
    {\fontsize{12pt}{14pt}\selectfont 
    IT Süsteemide Arendus \par}

    \vspace{2.5cm}

    % 3. Author (12pt)
    {\fontsize{12pt}{14pt}\selectfont Sergei Ivanov \par}

    \vspace{1cm}

    % 4. Title (16pt, BOLD, ALL CAPS)
    % Note: Subtitles should also follow the size/bold rules if part of the main title
    {\fontsize{16pt}{20pt}\selectfont \textbf{Lõputöö - essee ja tagasiside} \par}

    \vspace{1cm}

    % 5. Type of work (12pt)
    {\fontsize{12pt}{14pt}\selectfont Lõpuessee ja tagasiside \par}

    \vspace{3cm}

    % 6. Supervisor (12pt) - Position/Degree required
    \begin{flushright}
    {\fontsize{12pt}{14pt}\selectfont 
    Õppejõud: Jaak Vilo}
    \end{flushright}

    \vfill

    % 7. Location and Year (12pt)
    {\fontsize{12pt}{14pt}\selectfont Narva 2026 \par}

\end{titlepage}

\renewcommand{\abstractname}{Kokkuvõte}
\begin{abstract}
Käesolev aruanne analüüsib digipööret tarkvaraarenduse valdkonnas, 
keskendudes üleminekule kõrgema taseme raamistikelt süsteemprogrammeerimiskeeltele 
nagu Rust. Autor käsitleb kaasaegseid väljakutseid, sealhulgas "vibe-coding" ja
 tehisintellekti poolt genereeritud ebakvaliteetse koodi levikut, mis ohustab 
 süsteemide turvalisust ja hooldatavust. Töö raames hinnatakse kursuse 
 „Digitaalne maailmapilt“ rolli arendaja kompetentsi tõstmisel, 
 pakkudes võimalust tuletada meelde arvutiteaduse alustalasid ja 
 rakendada neid praktilistes projektides. Tulevikuvisioonis nähakse 
 tehisintellekti rolli pigem nõustaja kui iseseisva insenerina, rõhutades
  inimese rolli säilimist kriitilise tähtsusega süsteemide arhitektina ja vigade haldajana.
\end{abstract}

\vspace{0.5cm}

\noindent\textbf{Võtmesõnad:} Digipööre, tarkvaraarendus, Rust, vibe-coding, tehisintellekt, süsteemid, turvalisus, hooldatavus, digitaalne maailmapilt, arendaja kompetents, arvutiteaduse alused, praktilised projektid, tulevikuvisioon, nõustaja, insener, arhitekt, vigade haldamine

\renewcommand{\contentsname}{Sisukord}
\tableofcontents
\clearpage


